\section{Introduction}
\label{sec:intro}

\textbf{Domain.} Users judge an interactive service by its slowest responses.
When one user action fans out to many servers, a rare slow response on any
one of them delays the whole action, so the tail of each server (its slowest
1\%, the p99), not its average, sets the latency the user
sees~\cite{dean2013tail}. Many slow requests are not the application's
fault. The operating system delays them: the thread is ready to run but
waits for a CPU, network interrupts take over the core in the middle of a
request, the thread is moved to another CPU, or a write waits for the disk.
Li et al.~\cite{li2014tales} showed that such effects can make the tail of
simple Linux servers an order of magnitude worse than queueing theory
predicts.

eBPF lets us watch these events on a running system at low cost: small
programs, checked by the kernel's verifier, run at kernel tracepoints and
keep their results in kernel data structures called
maps~\cite{kernelbpf}. The usual eBPF tools, however, report totals: a
histogram of runqueue waits for a process, softirq time per CPU, counts per
second~\cite{bcc}. An operator who sees the p99 rise together with context
switches in the same minute still cannot tell whether the switches
\emph{caused} the slow requests, which requests they hit, or whether those
requests simply queued behind a burst of load.

\textbf{Problem.} For a service whose requests can be timed one by one, we
want to (i)~find the slow (tail) requests, (ii)~record for each one the
kernel events that delayed it, and (iii)~find which kernel signals really
separate slow requests from normal ones. Tracing may cost at most 5\% of
throughput, and a correlation must not be reported as a cause. Three things
make this hard. A request lasts only tens of microseconds, so even two
switches into the kernel per request are expensive. A request is often
delayed by something that happened to the requests \emph{ahead} of it in
the queue. And in a virtual machine, the host can pause the virtual CPUs
without the guest kernel noticing.

\textbf{Contributions.}
\begin{itemize}
\item \sys{}, an eBPF tool that keeps per-thread counters of preemption
wait, wakeup delay, blocked time and interrupt time (charged to the thread
that was interrupted) in a BPF map that user space can map into memory.
ghOSt's BPF scheduler shares per-task state through the same kind of
map~\cite{rhoden2022ghostbpf}; we use it to explain latency instead. The
application copies its own counters at the start and end of each request
with plain memory loads, with no system call or trap. This costs
\ovNsKvShm--\ovNsCpuShm{}\,ns of tracer CPU per request, 47--53$\times$ less
than uprobes (\S\ref{sec:overhead}).
\item A breakdown of each request's latency that includes delay
\emph{inherited} from requests queued ahead of it, checked by an accounting
identity (Eq.~\eqref{eq:identity}) that holds within 5\,$\mu$s for at least
\valWithin\% of requests in every run.
\item An evaluation over three workloads, three load levels and five
repetitions. It compares per-request kernel data with system-wide event
counts and with a queue-length baseline, on a workload the models never saw
and on fresh runs (H1), and checks the attribution by adding one stressor
at a time (H2).
\item Two by-products: a noise gate that detects hypervisor interference in
VM experiments, and a tested patch for bcc's \texttt{runqlat} that separates
waits after a wakeup from waits after a preemption.
\end{itemize}

\textbf{Hypotheses.} \textbf{H1}: kernel data aligned to each request's
time window separates slow requests from normal ones better than
system-wide event counts. \textbf{H2}: each controlled stressor raises the
delay component we predicted for it.
